% Results, sections 4.6 (jet mass: regression and sculpting) and 4.7 (X->bb and X->bs).
% Every number is a macro from experiments/FIGS/make_tables.py; CHECK comments mark claims to recheck.

\subsection{Jet mass: regression and sculpting}\label{sec:mass}

The jet-mass results use the JetClass-II test jets that have a generator-level match, so that the
generator-level soft-drop mass $\mgen$ is defined for resonant and QCD jets alike; a jet
without a match has no $\mgen$ and is left out. Every mass estimate is scored on the same
jets through its response $r$ and the resolution of equation~(\ref{eq:resolution}), and the binned
results use the reconstructed jet $p_{\mathrm{T}}$. The reference is the reconstructed soft-drop mass
$\msd$, which no training changes. The networks give two further estimates: the regressed mass
of the \Kres- and \Kcoarse-class models with the auxiliary mass regression, and an MLP regression on the
frozen features of every label set. Each model with the auxiliary mass regression has
\MlstMassNRunsWord{} pretraining runs; its tables give both values, and the text quotes the one less
favorable to the claim. The sculpting
results use the QCD jets of the JetClass-II test set and the output scores of the four levels of the
class hierarchy, of the two models with the auxiliary mass regression and of the three random partitions
with controlled flavor boundaries. The full output scores were extracted only for these models, so the
other random partitions are not part of the sculpting study.

The auxiliary mass regression gives a narrower mass response than the reconstructed soft-drop mass.
Over all test jets, the regressed mass of the \Kres-class model with the auxiliary mass regression % CHECK: MlstMassHeadResMassOverUncorr below 1 in both runs
reaches \MlstMassHeadResMassOverUncorrPr{} of $\msd$, that of the \Kcoarse-class % CHECK: MlstMassHeadCoarseMassOverUncorr not above MlstMassHeadResMassOverUncorr run by run, and within 2% of it (today 0.985 and 0.998)
model is as narrow, and both stay narrower in every bin of jet $p_{\mathrm{T}}$ and of % CHECK: every pT bin of F9b and every pT and m_gen bin of A17-res and A17-qcd has head resolution below groomed
$\mgen$ (figures~\ref{fig:f9a} and~\ref{fig:f9b}, and figures~\ref{fig:a17res} and~\ref{fig:a17qcd};
table~\ref{tab:massdetail}). The regressed mass also reaches % CHECK: MlstMassHeadResMassOverMlpRes below 1 in both runs
\MlstMassHeadResMassOverMlpResPr{} of an MLP regression on the frozen features of the \Kres-class
model pretrained without it, although that regression adds back the median of each jet's true class
and the auxiliary mass regression never sees the class. Because its target $\theta$ of
equation~(\ref{eq:masstarget}) is $\ln(\mgen/\msd)$, the regression learns from the jet's constituents how detector
response, pileup removal and grooming have moved the reconstructed mass away from the generator-level
one, which is the spread that $\msd$ carries. That correction depends on the constituents
rather than on which classes the labels separate, which is likely why the two label sets reach the
same width.

\begin{figure}
\centering
\paperfig{\linewidth}{F9a}
\caption{Distribution of the jet mass response $r = m_{\mathrm{est}}/\mgen$ of equation~(\ref{eq:resolution}),
the ratio of a jet's mass to its generator-level mass, for (a) resonant jets and (b) QCD jets of the
JetClass-II test set with a generator-level match: the reconstructed soft-drop mass (dashed reference) and the
regressed mass of the \Kres- and \Kcoarse-class models with the auxiliary mass regression, each the mean over
its \MlstMassNRunsWord{} pretraining runs. The lower panels give the ratio of each distribution to that of the
soft-drop mass. The horizontal axis of (b) holds the soft-drop mass response of \MlstMassRespAxesShareQcd{}
of the QCD jets. The resolution and the median response are in table~\ref{tab:massdetail}.}
\label{fig:f9a}
\end{figure}

\begin{figure}
\centering
\paperfig{\linewidth}{F9b}
\caption{Median jet mass response (a) and relative mass resolution $\sigma_r$ of
equation~(\ref{eq:resolution}) (b) against the reconstructed jet $p_{\mathrm{T}}$, for all JetClass-II test
jets with a generator-level match: the reconstructed soft-drop mass and the regressed mass of the \Kres- and
\Kcoarse-class models with the auxiliary mass regression, each the mean over its \MlstMassNRunsWord{}
pretraining runs. Below are the ratio of each median to that of the soft-drop mass and the root-square
difference of each resolution and that of the soft-drop mass, as in ATLAS~\cite{atlasregression}.}
\label{fig:f9b}
\end{figure}

The auxiliary mass regression also puts the jet mass into the frozen features themselves. An MLP
regression on the frozen features of the \Kres-class model with the auxiliary mass regression % CHECK: MlstCmpMlpResMassOverRes below 1 in both runs (today 0.7186 and 0.7253)
reaches \MlstCmpMlpResMassOverResPr{} of the same regression on the \Kres-class model's features, % CHECK: MlstCmpMlpCoarseMassOverCoarse below 1 in both runs; the two regressed-feature resolutions agree (today 0.1583, 0.1588 against 0.1591, 0.1586) and the 17-class model starts narrower (0.2062, 0.2118 against 0.2214, 0.2187)
and the \Kcoarse-class pair gains a little less, \MlstCmpMlpCoarseMassOverCoarsePr{} of the \Kcoarse-class
model's, because that model starts narrower and both reach the same resolution (table~\ref{tab:mass}). Without the auxiliary target the
label sets differ far less (\Kcoarse- over \Kfull-class model, \MlstCmpMlpCoarseOverFullPr). The % CHECK: MlstCmpMlpCoarseOverFull much closer to 1 than MlstCmpMlpResMassOverRes
reason lies in the training data: the reweighting flattens every reweighting category in
$p_{\mathrm{T}}$ and $\msd$, so a jet's mass says little about its class, and no
classification label rewards features that keep the mass more precisely than the decay topology
implies. The regression term does. The row of the randomly initialized network in
table~\ref{tab:mass} shows how much of the mass a network that never trained already exposes, and % CHECK: MlstMassMlpInit above MlstMassMlpFull (pretrained features carry more mass than the random network)
every pretrained model's features carry more.

\begin{table}[t]
\centering
\caption{Jet-mass resolution and median response on JetClass-II, compared between the soft-drop mass, the mass output of the pretraining head and a multilayer perceptron that regresses the mass from the frozen features. The response $r = m/m_{\mathrm{gen}}$ is the estimated mass over the generator-level soft-drop mass, and the resolution is half the range between the 15.9\% and 84.1\% quantiles of $r$ divided by its median, over all test jets with a generator-level match. The perceptron's target is centered on the training median of the jet's true class, which the other two estimates are not. Values are the mean $\pm$ standard deviation over \MlstNRunsWord{} pretraining runs per label set, \MlstMassNRunsWord{} per network with an added mass output and three initializations of the random network.}
\label{tab:mass}
\begin{tabular*}{\linewidth}{@{\extracolsep{\fill}}lcc@{}}
\toprule
 & resolution & median response \\
\midrule
soft-drop mass & 0.2448 & 0.9294 \\
\midrule
\multicolumn{3}{@{}l}{\textit{mass output of the pretraining head}} \\
\Kres{} classes + mass & 0.1467\,$\pm$\,0.0006 & 0.9992\,$\pm$\,0.0055 \\
\Kcoarse{} classes + mass & 0.1455\,$\pm$\,0.0008 & 1.0052\,$\pm$\,0.0002 \\
\midrule
\multicolumn{3}{@{}l}{\textit{multilayer-perceptron probe on the frozen features}} \\
\Kfull{} classes & 0.2200\,$\pm$\,0.0014 & 0.9912\,$\pm$\,0.0047 \\
\Kres{} classes & 0.2195\,$\pm$\,0.0017 & 0.9876\,$\pm$\,0.0014 \\
\Kmid{} classes & 0.2158\,$\pm$\,0.0029 & 0.9908\,$\pm$\,0.0049 \\
\Kcoarse{} classes & 0.2098\,$\pm$\,0.0031 & 0.9908\,$\pm$\,0.0041 \\
\Kres{} classes + mass & 0.1589\,$\pm$\,0.0003 & 1.0049\,$\pm$\,0.0015 \\
\Kcoarse{} classes + mass & 0.1585\,$\pm$\,0.0003 & 1.0077\,$\pm$\,0.0022 \\
random network & 0.2421\,$\pm$\,0.0018 & 0.9833\,$\pm$\,0.0049 \\
\bottomrule
\end{tabular*}
\end{table}


The auxiliary mass regression costs the classification of the pretraining classes little. A linear
probe on the frozen features of the \Kres-class model with the auxiliary mass regression classifies % CHECK: MlstCmpJcTwoResMassOverRes today 1.022 and 1.016; its PrMax reads "equal within" under the Methods two-run rule or "up to ... more errors", and both fit "costs little"
JetClass-II about as well as one on the \Kres-class model (ratio of their $1-\AUC$:
\MlstCmpJcTwoResMassOverResPrMax{}; table~\ref{tab:probes}). At \Kcoarse{} classes, where the mass term takes a larger share of the % CHECK: MlstMassLossShareCoarseMass above MlstMassLossShareResMass (today 28.7 vs 13.8 per cent)
loss (section~\ref{sec:labelsets}), the two runs move in opposite directions, so the larger share % CHECK: MlstCmpJcTwoCoarseMassOverCoarse straddles 1 (today 1.054 and 0.977)
does not show a larger cost. A likely reason is that the frozen features have room for the mass
beside the class information, so the regression adds what the features carry without displacing
what the labels ask for. With \MlstMassNRunsWord{} runs per model the data can only fail to show a
cost that grows with the loss share; they cannot exclude a modest one. An MLP probe on the natural-composition
test set reads the same at both label sets (ratio of its $1-\AUC$: \MlstMlpNatJcTwoResMassOverResPr{} at \Kres{}
classes and \MlstMlpNatJcTwoCoarseMassOverCoarsePr{} at \Kcoarse{} classes; table~\ref{tab:mlp}), so this is not
a limit of what a linear classifier can read.
% CHECK: MlstMlpNatJcTwoResMassOverResPr and MlstMlpNatJcTwoCoarseMassOverCoarsePr read "equal within" (recomputed from fits_v3_mlp2: raw 1.016, 1.009 and 1.026, 1.019; hierarchy-level SD 1.3%).

Fine-tuning on top tagging shows no change from the auxiliary mass regression beyond the uncertainty of
$R_{30}$ at either label set. With it, the \Kres-class runs move $R_{30}$ in opposite % CHECK: MlstFtTopRthirtyResMassOverRes straddles 1 (raw R30 today 1.12 and 0.89)
directions (\MlstFtTopRthirtyResMassOverResPr), and the \Kcoarse-class runs give % CHECK: MlstFtTopRthirtyCoarseMassOverCoarsePr reads "equal within" under the Methods two-run rule (raw R30 today 0.90 and 0.92); if it states a loss, it must also exceed MlstBootTestMaxRthirtyFt before the text may call it one
\MlstFtTopRthirtyCoarseMassOverCoarsePr{}, within the counting uncertainty of a fine-tuned $R_{30}$ ratio
between label sets, up to \MlstBootTestMaxRthirtyFt\% if the two counts are taken as independent (table~\ref{tab:ftrej};
appendix~\ref{sec:app-boot}). As on the pretraining
classes, the larger loss share at \Kcoarse{} classes shows no larger cost; with \MlstMassNRunsWord{} runs
per model a modest one cannot be excluded, and a pair of models with the same loss share at both label
sets would test it.

% PI 2026-11-09: the main text reports exactly the sculpting differences beyond 3 sigma (tab:sculptboot,
% gallery_placement.F10B_STATED, tested against sculpting_boot.json) and nothing else; figures, tables and the
% floor are in appendix sec:app-sculpt.
A cut on the pretraining head of the \Kmid- or the \Kcoarse-class model sculpts the QCD mass distribution
more than a cut on the \Kfull-class head, each at one working point, and the auxiliary mass regression reduces
the sculpting by the \Kres-class head. These are the only differences between pretraining heads that exceed
three times their uncertainty (table~\ref{tab:sculptboot}): the \Kcoarse-class head at a QCD efficiency of
\SculptEffLow, the \Kmid-class head at \SculptEffHigh{} and the \Kres-class head with the auxiliary mass
regression at \SculptEffHigh. Every other difference lies within its uncertainty, among them the auxiliary mass
regression at \Kcoarse{} classes and the three random partitions with controlled flavor boundaries against
the \Kfull-class model (appendix~\ref{sec:app-sculpt}).

% BEGIN XBX
\subsection{$X\to b\bar b$ and $X\to bs$ resonance tagging}\label{sec:xbx}

The resonance benchmark repeats Sophon's~\cite[section~III]{sophon} at every label set, on samples of % CHECK: pending XbxMx and the spin and production of the staged samples (Sophon: spin-0)
a resonance $X$ decaying to $b\bar b$ or to $bs$, with the background, mass window and metrics of
section~\ref{sec:evaluation}; figure~\ref{fig:f11} gives the resonance mass, the mass window and the
numbers of jets. The transfer heads are trained on QCD jets from files that no pretraining run read,
drawn to match the signal in $p_{\mathrm{T}}$ and $\msd$, so no jet seen in pretraining enters
the benchmark, and every model is scored on the same signal and QCD jets. Figure~\ref{fig:f11} draws
$\epsS/\sqrt{\epsB}$ of equation~(\ref{eq:signif}), the reach of a selection without event
counts, where Sophon plots the significance $Z$ itself~\cite[figure~2]{sophon}; our curves rank the label
sets on one background.

Transfer from frozen features tags $X\to bs$ better than a network trained from scratch for the task, % CHECK: pending XbsRthirtyFullOverScratch above 1 at tight working points (tab:xbxrej; F11 draws no network trained from scratch)
at tight working points. Measured by $R_{30}$, a transfer head on the frozen features of the
\Kfull-class model gives \XbsRthirtyFullOverScratchPr{} than the network trained from scratch on the
same jets (table~\ref{tab:xbxrej}), and the gain depends on the label set: the \Kcoarse-class % CHECK: pending XbsRthirtyCoarseOverFull below 1
features give \XbsRthirtyCoarseOverFullPr{} than the \Kfull-class ones. The $X\to bq$ output node of the % CHECK: pending XbsRthirtyNodeOverTransferFull below 1 (closest node weaker than the transfer head)
\Kfull-class model, scored against QCD, tags $X\to bs$ worse than the transfer head, so the gain sits
in the frozen features rather than in an output the network already had (table~\ref{tab:xbxrej}).
An $X\to bs$ jet is told from QCD by two prongs of which exactly one carries a $b$ quark. The
\Kfull-class labels separate two-prong decays by the flavor of each prong, while the \Kcoarse-class
labels merge them, so the coarser features likely keep less of the single-$b$ information this signal
needs.

Sophon's direct discriminant needs no new training but exists only where $X\to b\bar b$ is a class of
its own, at \Kfull{} and \Kres{} classes (table~\ref{tab:usecase}). On the \Kfull-class model it % CHECK: pending XbbRthirtyDirectOverTransferFull within 2 SD of 1; if not, state which is better
reaches the rejection of a transfer head trained on the same frozen features
(\XbbRthirtyDirectOverTransferFullPr; figure~\ref{fig:f11}(a)), so the output nodes already carry what
the head learns for this signal. The transfer head compares every label set, and the \Kcoarse-class % CHECK: pending XbbRthirtyTransferCoarseOverFull below 1
features lose rejection against the \Kfull-class ones (\XbbRthirtyTransferCoarseOverFullPr). The
\Kcoarse-class labels merge $X\to b\bar b$ with two-prong decays of other quark flavors, likely the
reason, since that merge is also where a linear probe loses most in separating $X\to b\bar b$ from
$X\to c\bar c$ (section~\ref{sec:jetclasstwo}).

\begin{figure}
\centering
\paperfig{\linewidth}{F11}
\caption{Resonance tagging of (a) $X\to b\bar b$ and (b) $X\to bs$ against QCD jets: $\epsS/\sqrt{\epsB}$ of
equation~(\ref{eq:signif}) against the QCD efficiency $\epsB$, for Sophon's direct discriminant of the
\Kfull-class model and a transfer head on the frozen features of each label set, each the mean over runs; a
curve ends where fewer than $\nmin$ QCD jets pass. Jets lie in the soft-drop mass window
\XbxWindowLo--\XbxWindowHi~GeV around a resonance mass of \XbxMx~GeV, and every model is scored on the same
\XbbNEvalSignal{} $X\to b\bar b$, \XbsNEvalSignal{} $X\to bs$ and \XbsNEvalQcd{} QCD jets. The curves are
pending the $X\to b\bar b$ and $X\to bs$ samples.}
\label{fig:f11}
\end{figure}
% END XBX
